\section{Related Literature}
\label{sec:literature}

\subsection{Tournament Labor Markets}
\label{sec:lit_tournaments}

The setting studied here is a tournament labor market in the sense of \citet{Lazear1981_tournaments}: compensation and advancement depend on rank rather than on absolute output. Rank-order tournaments economize on measurement when absolute performance is costly to observe, and \citet{Rosen1986_elimination_tournaments} shows that in elimination structures the prize spread must widen at later rounds to maintain effort, because the option value of continuing shrinks as the field narrows. \citet{Prendergast1999_incentives_firms} surveys the resulting incentive literature. Empirical tests have relied heavily on sport, where rank, prize, and effort are all measured: \citet{Ehrenberg1990_tournaments_golf} show that prize spread affects scoring in professional golf, \citet{Sunde2009_heterogeneity_tennis} finds that heterogeneity between competitors dampens effort in tennis, and \citet{Genakos2012_interim_rank} show that interim rank shapes risk-taking in dynamic tournaments. \citet{Kahn2000_sports_laboratory} sets out the general case for sport as a labor-market laboratory.

This paper takes the tournament structure as given and asks a question the tournament literature has not addressed: what happens when \textit{access} to the tournament is itself randomly assigned. Standard tournament models treat the set of competitors as fixed and study effort within a contest. Professional tennis instead runs a sequence of contests in which yesterday's rank determines today's entry, so a shock to rank propagates as a shock to future participation. The object of interest is therefore the return to an entry slot rather than the elasticity of effort with respect to the prize spread.

\subsection{Randomized Access Shocks}
\label{sec:lit_access}

The closest analog is \citet{Pallais2014_inefficient_hiring}, who randomly assigns first job opportunities to inexperienced workers on an online labor market and finds large, persistent employment gains. The mechanism she identifies is informational: the randomly granted job generates a public rating that subsequent employers use, so the market's own record-keeping converts a one-time assignment into a durable advantage. The tennis ranking plays the same role as the oDesk rating, and the present setting shares the essential design feature of literal randomization.

Four features of this setting let the analysis go beyond what the oDesk design supports, and they define the paper's contribution relative to \citet{Pallais2014_inefficient_hiring}. First, outcomes are observed at five horizons from 4 to 52 weeks rather than at a single follow-up, so the trajectory of the effect is identified rather than assumed. Second, the market reports two distinct measures of a worker: ranking points, which reward participation without adjusting for opponent quality, and Elo, which adjusts for it. Their divergence separates an access channel from an ability channel, a decomposition that a single reputational score cannot deliver. Third, the effect is observed to decay, and the decay rate is informative about the market's memory: the 52-week rolling ranking window mechanically retires the initial shock, whereas an oDesk rating persists indefinitely. Fourth, some workers receive the same randomized shock more than once, which identifies how the return to an opportunity varies with the number already received.

Related evidence on access shocks comes from \citet{Abel2020_reference_letters}, who randomize reference letters for South African job seekers, and \citet{Startz2021_interviewer_effects}, who exploits random interviewer assignment. In science, \citet{Azoulay2014_matthew_effect} document cumulative advantage following early recognition, \citet{Wang2019_early_career_setback} find that narrowly missing early funding predicts stronger later performance, and \citet{Azoulay2010_superstar_extinction} show that losing access to a star collaborator reduces output. Threshold-crossing designs in education find long-run earnings gains from marginal admission \citep{Zimmerman2014_college_marginal, Hoekstra2009_flagship_earnings, Goodman2017_college_access}.

\subsection{Initial Conditions and Path Dependence}
\label{sec:lit_initial}

A large literature establishes that early-career conditions have durable effects: graduating in a recession depresses earnings for a decade \citep{Oreopoulos2012_graduating_recession}, initial placement shapes the subsequent research careers of economists \citep{Oyer2006_initial_placement}, MBA cohort timing determines entry into investment banking \citep{Oyer2008_investment_banker}, and childhood location shapes intergenerational mobility \citep{Chetty2014_where_is_land}. \citet{Malmendier2011_depression_babies} show that early macroeconomic experience alters risk preferences decades later, which is the sharpest existing statement of the claim that a first experience is qualitatively different from a later one. The common mechanism is cumulative advantage \citep{Merton1968_matthew_effect, Rosen1981_superstars}.

This paper differs from that literature in the nature of the shock and in the market that transmits it. The shocks studied there are shocks to the environment a career starts in; the shock here occurs within an ongoing career and is randomly assigned by an institutional rule. The persistence results also differ, and the difference is informative: where recession-cohort effects run for a decade, the access gains documented here dissipate within a year. Section~\ref{sec:conclusion} argues that this contrast identifies market transparency, and specifically the presence of a rolling performance window, as the moderator governing how long an access shock survives.

\subsection{Sports as an Empirical Laboratory}
\label{sec:lit_sports}

\citet{Gauriot2019_fooled_randomness} use narrowly decided tennis matches to show that ranking systems reward luck as though it were skill, which motivates the separation of ranking points from Elo pursued here. \citet{GonzalezDiaz2012_performing_best} study performance under pressure and \citet{Berger2011_can_losing} examine responses to interim deficits in basketball.

Most directly, \citet{Maity2025_early_career_sports} study lucky losers across several professional sports and report that early setbacks correlate with stronger subsequent performance. Their comparison is between lucky losers and other qualifying losers without a design that accounts for how lucky losers are selected, so the estimates combine the effect of entry with the selection that produced it. The Grand Slam lottery provides the missing randomization, and Section~\ref{sec:mechanisms_reconciling} shows that the correlation they document is consistent with an access channel rather than a performance channel.
